\documentclass[UTF8,fontset=none,11pt,a4paper]{ctexart}
\usepackage{ipara-handbook}
\handbooksetup{DS-09 查找与有序索引}{408 · 数据结构 DS-09}{Order · Search · Balance · Update}
\begin{document}
\handbooktitle{查找与有序索引}{维护多少秩序，才能换取多少查询能力}{408 · 数据结构 Topic DS-09}

\begin{corebox}{母问题：查询加速从哪里来，又由谁付费？}
有序索引的生成链是
\[
\text{Order/Index}\to\text{Search Reduction}\to\text{Update Repair}\to\text{Height Invariant}\to\text{Cost Tradeoff}.
\]
顺序查找几乎不预处理；折半查找要求随机访问和有序区间；BST 把次序放进树路径；平衡树用旋转或分裂把高度约束在对数级。查询越强，通常越要在插入、删除、空间或维护时间上付费。
\end{corebox}
\begin{methodbox}{本册定位与 Owner 边界}
\begin{itemize}
\item \kw{Owns：}顺序/分块/折半查找、BST、AVL/红黑树的有序性与高度维护、B 树的查找与插删、B+ 树的查找结构。
\item \kw{Uses：}Atlas Foundation 成本向量；DS-B02 负责按 workload 比较有序索引与 Hash。
\item \kw{Stops：}Hash 映射归 DS10；磁盘块 I/O 的外部排序归 DS12；具体工程库实现只作 Extension。
\end{itemize}
\end{methodbox}
\tableofcontents\newpage

\section{查找契约先固定}
“查找”要先回答返回的是存在性、位置、前驱、首个重复 key 还是范围边界。折半查找依赖一个明确的区间不变量，例如目标若存在必在 $[low,high]$ 中；每轮比较后必须缩小区间，空区间才是失败证明。重复 key 若不固定取左界/右界，算法结果虽都可能“找到”却不满足接口。

\subsection{顺序、分块与折半：三种秩序预算}
顺序查找不要求数据有序，逐项排除候选。分块查找把表分成若干块，维护“块间有序、块内可无序”的两级结构：先在索引表定位可能的块，再在块内顺序查找。其平均成本可写成
\[
ASL\approx ASL_{index}+ASL_{block},
\]
因此块数增加会让索引阶段变长、块内阶段变短；不能只优化其中一项。折半查找则要求整体有序且支持随机访问，每轮用中点把候选区间近似减半。

三者逐步增加维护秩序以缩小查询候选：无索引、稀疏块索引、全局有序随机访问。若数据频繁插入，维持整体连续有序的移动成本可能抵消查询收益。

\section{BST：次序下沉到路径}
BST 对每个节点保持左子树 key 小于节点、右子树 key 大于节点（重复值策略需显式约定）。查找、插入、删除都沿一条根叶路径；复杂度是 $O(h)$，不是无条件 $O(\log n)$。退化链的高度为 $n$，所以平衡条件不是装饰而是复杂度证明的前提。
\begin{examplebox}{删除的三种结构分支}
叶节点直接删除；单孩子用孩子替代节点；双孩子从中序后继/前驱取替代 key，再在其原位置删除。每一步都要恢复局部 BST 有序性，不能只交换值而忘记释放或接回子树。
\end{examplebox}

\section{AVL 与红黑树：都控制高度，但修复契约不同}
AVL 约束任一节点左右子树高度差不超过 1。插入后从修改点向根寻找第一个失衡祖先，根据较高子树方向形成 LL/RR/LR/RL；旋转后必须同时恢复 BST 中序有序和高度约束。

红黑树把严格高度平衡放宽为颜色约束：根与空叶为黑；红节点的孩子不能为红；从任一节点到其后代空叶的每条路径具有相同黑高。由此最长根叶路径不超过最短路径两倍，树高仍为 $O(\log n)$。插入冲突按叔节点颜色和结构方向决定变色或旋转；删除若移走黑色会产生黑高亏损，修复目标是把亏损上移、吸收或消除。

\begin{warnbox}{AVL 与红黑树不能共享旋转口诀}
AVL 的触发量是高度差，红黑树的触发量是颜色与黑高。二者旋转都保持 BST 次序，但合法状态不同；只看到“失衡就旋转”无法判断变色与停止条件。
\end{warnbox}

\section{B 树：用节点扇出降低外存访问高度}
设 $m$ 阶 B 树，每个节点至多有 $m$ 个孩子、至多 $m-1$ 个关键字；除根外的非叶节点至少有 $\lceil m/2\rceil$ 个孩子，所有叶节点位于同一层。节点内关键字有序并划分孩子区间。查找先在节点内确定区间，再沿一个孩子下降；昂贵的是访问多少个节点/块。

插入总在叶层落位。节点溢出时把中间关键字上升父节点，左右关键字分成两个合法节点；父节点若再溢出就继续向上，根分裂会使树高增加 1。删除后若节点低于下限，先向兄弟借一个关键字并经父节点重分配；不能借才与兄弟及父分隔关键字合并。若根最终无关键字且只有一个孩子，用该孩子替代根，树高减 1。

统一不变量是：节点占用合法、节点内有序、孩子区间正确、所有叶同层。分裂和合并是在保持这些不变量下控制高度。

\section{B+ 树：内部节点导航，叶层拥有完整检索序列}
B+ 树把记录或记录指针集中在叶节点；非叶节点只保存导航分隔 key，关键字可能在内部和叶层重复。所有叶节点按 key 顺序链接，因此等值查找沿根到叶完成，范围查找定位起点后沿叶链扫描。

\begin{center}\small
\begin{tabularx}{\linewidth}{L{2.0cm}C{.35cm}L{2.0cm}YY}\toprule
概念 A&$\neq$&概念 B&真正判据与题面信号&混淆后果\\\midrule
B 树&$\neq$&B+ 树&B 树各层可命中记录；B+ 树内部只导航、叶链承载全体 key&误判成功层次和范围扫描路径\\
树的阶&$\neq$&节点关键字数&阶限制最大孩子数，关键字上限通常少 1&占用上下限算错\\
分裂&$\neq$&旋转借位&分裂制造新节点；借位在相邻节点和父分隔 key 间重分配&高度变化判断错\\\bottomrule
\end{tabularx}\end{center}

\section{代码与测试证据}
完整实现位于 \codepath{code/ds09_ordered_index.hpp}，测试位于 \codepath{code/ds09_ordered_index_test.cpp}。本册用边界明确的 BST 作为可执行核心，平衡树和 B/B+ 的分裂规则保留为机制扩展。
\lstinputlisting[language=C++,firstline=12,lastline=104,firstnumber=12,numbers=left,caption={BST 查找、插入、删除与有序输出}]{30_408/10_数据结构/09_查找与有序索引/code/ds09_ordered_index.hpp}
\lstinputlisting[language=C++,firstline=8,lastline=44,firstnumber=8,numbers=left,caption={空树、重复 key 与删除分支测试}]{30_408/10_数据结构/09_查找与有序索引/code/ds09_ordered_index_test.cpp}
\begin{lstlisting}[language=bash]
clang++ -std=c++17 -Wall -Wextra -Werror -pedantic code/ds09_ordered_index_test.cpp -o /tmp/ds09_test
/tmp/ds09_test
\end{lstlisting}

\section{复杂度与概念边界}
\begin{center}\small
\begin{tabularx}{\linewidth}{L{2.5cm}C{1.8cm}C{1.8cm}YY}\toprule
结构/操作 & 平均或保证 & 最坏 & 依赖的不变量 & 维护代价\\\midrule
顺序查找 & $O(n)$ & $O(n)$ & 无序也可 & 几乎无预处理\\
折半查找 & $O(\log n)$ & $O(\log n)$ & 有序+随机访问 & 插入移动\\
BST & $O(h)$ & $O(n)$ & 局部有序 & 可能退化\\
平衡树 & $O(\log n)$ & $O(\log n)$ & 高度/颜色约束 & 旋转/分裂\\\bottomrule\end{tabularx}\end{center}

\section{做题控制协议}
先写返回契约和区间/树不变量，再看输入是否有序、是否支持随机访问、更新是否频繁。看到 Hash 适合精确等值但不适合范围查询；看到动态有序更新，才考虑平衡索引。每个复杂度都把高度或有序前提写出来。
\begin{corebox}{压缩复原}
五问：查询要返回什么？候选区间怎样缩小？秩序存在哪里？更新破坏哪个不变量？维护成本由谁承担？
\end{corebox}
\end{document}
